\documentclass[11pt]{article}
\usepackage[margin=1in]{geometry}
\usepackage{amsmath,amssymb}
\usepackage{enumitem}
\usepackage{xcolor}
\usepackage[colorlinks=true,linkcolor=blue!50!black,urlcolor=blue!50!black]{hyperref}
\setlist{nosep,leftmargin=1.4em}
\newcommand{\stated}{\textcolor{green!40!black}{\textbf{Stated by the authors:}}}
\newcommand{\ours}{\textcolor{orange!80!black}{\textbf{Our reading:}}}
\usepackage{xurl}
\newcommand{\file}[1]{\url{#1}}
\sloppy
\setlength{\parindent}{0pt}\setlength{\parskip}{4pt}

\title{Beyond Homogeneous Agents\\[2pt]\large Heterogeneity, incentives, adversaries and collisions}
\author{CS6007 project notes}
\date{October 2026}
\begin{document}
\maketitle
\textbf{What this folder is about.} When agents are identical and cooperative, a network
of $N$ agents is ``one agent with $N$ times the data'', which we proved for our setting. The
papers here study what makes multi-agent bandits \emph{genuinely different}: agents that differ
(heterogeneous or clustered models), agents with their own interests (incentives), agents that
lie (malicious recommendations), and agents that compete for the same arms (collisions). This is
the folder to read for new research directions; E29 (heterogeneous slopes) belongs here.
\tableofcontents
\bigskip
\section{Shi \& Shen (AAAI 2021): federated multi-armed bandits}
\textbf{C. Shi, C. Shen. Federated Multi-Armed Bandits. AAAI 2021 (arXiv:2101.12204).}\\ File: \file{2101.12204v2.pdf}
\paragraph{Setting.}
\begin{itemize}\item Clients have \emph{heterogeneous local} bandit models; the goal is the \emph{global} model. In the approximate model, local models are random realisations of the global one, so the global model is uncertain even if local models were known.\end{itemize}
\paragraph{What they did.}
\begin{itemize}\item \textbf{Fed2-UCB}: a ``double UCB'' accounting for both arm-sampling and client-sampling uncertainty, admitting clients gradually.\item \textbf{Fed1-UCB} for an exact model where the global model is the average of local ones.\end{itemize}
\paragraph{Main results.}
\begin{itemize}\item $O(\log T)$ regret for the approximate model while controlling communication; admitting clients gradually is critical for that rate.\item An order-optimal exact-model result that trades convergence time against the number of clients.\end{itemize}
\paragraph{Why it matters for our project.}
\begin{itemize}\item A clean example of a \emph{new source of uncertainty} that exists only with many agents (client sampling): the kind of contribution we are looking for.\end{itemize}
\paragraph{Open problems.}
\begin{itemize}\item \stated{} The risk when synchronisation breaks is noted as ``an interesting research problem for future investigation''.\item \ours{} Multi-armed only; contextual or structured (e.g.\ single-index) versions with heterogeneous clients are open.\end{itemize}
\section{Ghosh, Chung, Yin \& Ramchandran (NeurIPS 2020): clustered federated learning (IFCA)}
\textbf{A. Ghosh, J. Chung, D. Yin, K. Ramchandran. An Efficient Framework for Clustered Federated Learning. NeurIPS 2020 (arXiv:2006.04088).}\\ File: \file{2006.04088v2.pdf}
\paragraph{Setting.}
\begin{itemize}\item Federated learning (supervised, not bandits) where users belong to unknown clusters, each with its own model.\end{itemize}
\paragraph{What they did.}
\begin{itemize}\item \textbf{IFCA}: alternately estimate each user's cluster identity and update each cluster's model by gradient descent.\end{itemize}
\paragraph{Main results.}
\begin{itemize}\item Convergence guarantees for linear models with squared loss and for strongly convex smooth losses, given good initialisation; a discussion of optimal statistical error.\item In practice, random initialisation with restarts works, including beyond the convex regime.\end{itemize}
\paragraph{Why it matters for our project.}
\begin{itemize}\item Written by our instructor. The natural template for \emph{clustered single-index bandits}: agents with different $\theta^*_i$, where naive pooling hurts (our E12).\end{itemize}
\paragraph{Open problems.}
\begin{itemize}\item \stated{} Weakly convex and non-convex losses.\item \stated{} Stochastic gradients on the workers.\item \stated{} Only a small subset of devices participating.\item \ours{} Not a bandit paper: clustering while \emph{also} exploring (online, with partial feedback) is open.\end{itemize}
\section{Wei, Li, Xu \& Wang (2023): incentivised communication for federated bandits}
\textbf{Z. Wei, C. Li, H. Xu, H. Wang. Incentivized Communication for Federated Bandits. arXiv:2309.11702, 2023.}\\ File: \file{2309.11702v2.pdf}
\paragraph{Setting.}
\begin{itemize}\item Federated linear contextual bandits where clients are \emph{self-interested} and share data only if it benefits them; the server must pay incentives.\end{itemize}
\paragraph{What they did.}
\begin{itemize}\item \textbf{Inc-FedUCB}: a general framework with two implementations, data incentives and monetary incentives.\end{itemize}
\paragraph{Main results.}
\begin{itemize}\item Customisable near-optimal regret with guarantees on both communication and payment costs; experiments across environments.\end{itemize}
\paragraph{Why it matters for our project.}
\begin{itemize}\item A multi-agent question with no single-agent counterpart. Exploration in our Phase 1 is a pure cost to the explorer, so the same free-riding issue applies.\end{itemize}
\paragraph{Open problems.}
\begin{itemize}\item \stated{} Clients currently report their costs truthfully; with strategic misreporting a \textbf{truthful incentive mechanism} is needed.\end{itemize}
\section{Vial, Shakkottai \& Srikant (2021): robust multi-agent bandits}
\textbf{D. Vial, S. Shakkottai, R. Srikant. Robust Multi-Agent Multi-Armed Bandits. arXiv:2007.03812, 2020 (v3 2021).}\\ File: \file{2007.03812v3.pdf}
\paragraph{Setting.}
\begin{itemize}\item $n$ honest and $m$ malicious agents; honest agents recommend their best-arm estimates, malicious ones recommend arbitrary arms.\end{itemize}
\paragraph{What they did.}
\begin{itemize}\item Shows existing collaborative algorithms lose all benefit with even \emph{one} malicious agent.\item Proposes \textbf{blocking}: honest agents learn who is malicious and reduce communication with them.\end{itemize}
\paragraph{Main results.}
\begin{itemize}\item With blocking, collaboration again reduces regret below the single-agent baseline, assuming $m$ is small compared with $K$ and without assumptions on malicious behaviour.\end{itemize}
\paragraph{Why it matters for our project.}
\begin{itemize}\item Shows recommendation-based protocols (like Chawla et al.) are fragile; any relay of statistics in our protocol has the same vulnerability.\end{itemize}
\paragraph{Open problems.}
\begin{itemize}\item \ours{} Multi-armed only; robustness for contextual or single-index bandits (where malicious agents could send poisoned statistics) is open.\item \ours{} Graph-structured versions exist in later work by the same authors.\end{itemize}
\section{Proutière \& Wang (2019): optimal multiplayer bandits with collisions (DPE)}
\textbf{A. Proutière, P.-A. Wang. An Optimal Algorithm for Multiplayer Multi-Armed Bandits. arXiv:1909.13079, 2019.}\\ File: \file{1909.13079v2.pdf}
\paragraph{Setting.}
\begin{itemize}\item $M$ players choose among $K$ arms; if several players pick the same arm they \emph{collide} and get nothing; players observe collisions. Not all players know $M$.\end{itemize}
\paragraph{What they did.}
\begin{itemize}\item \textbf{DPE} (Decentralized Parsimonious Exploration): a leader--follower algorithm in which players communicate only through deliberate collisions.\end{itemize}
\paragraph{Main results.}
\begin{itemize}\item Regret asymptotically matching the optimal \emph{centralised} algorithm, better than SIC-MMAB.\item The expected number of rounds used for communication is finite.\end{itemize}
\paragraph{Why it matters for our project.}
\begin{itemize}\item A completely different multi-agent problem (competition for arms, communication through actions). Its leader-announced commit time is a template for our Phase-1 agreement problem.\end{itemize}
\paragraph{Open problems.}
\begin{itemize}\item \ours{} Players share the same arm means; player-dependent rewards and the no-sensing (collisions not observed) variant need different ideas.\end{itemize}
\section{Open problems across this folder}
\begin{itemize}\item \textbf{Contextual or structured versions} of every idea here are mostly open: heterogeneous clients, clustering, incentives, robustness, all studied for MAB or linear models only.\item \textbf{Clustered online learning} (IFCA while exploring) is open, and fits single-index bandits with different $\theta^*_i$.\item \textbf{Heterogeneous identifiability} (our E29): agents whose own slopes $\mu_i$ differ in sign or vanish can identify $\theta^*$ only together. None of these papers covers it, because none uses an index model.\end{itemize}
\end{document}